\begin{annexitem}{annex:p09-timed-rescaling}{thm:p09-timed-rescaling}{Rescaling the clocks and all constants}
\begin{proof}
Let \(T_q(v)=qv\). It is a bijection of nonnegative valuations with inverse
\(T_{1/q}\), and
\(T_q(v+d)=T_q(v)+qd\), \(T_q(v[r])=T_q(v)[r]\).
For every positive \(q\), an atomic unary or difference comparison holds
iff its multiplied comparison holds: multiplication by \(q\) preserves
both weak and strict order and equality. Conjunction preserves this
equivalence. Therefore a delay preserving the invariant at all
\(t\in[0,d]\) corresponds to one preserving it at all
\(qt\in[0,qd]\); the latter interval is exactly the image of the former.
Guard and destination-invariant satisfaction correspond too. Induct on the
finite edge sequence to obtain the stated run bijection. Conversely divide
the values and delays by \(q\) at every step. A sum of nonnegative delays
is unbounded iff its positive constant multiple is unbounded, giving the
divergence statement without any acceptance-condition inference. A common
multiple of the finitely many denominators makes every constant integral.
\end{proof}
\end{annexitem}

\begin{annexitem}{annex:p09-timed-region-properties}{thm:p09-timed-region-properties}{The explicit cell partition and event-time matching}
\begin{proof}
The descriptor consists of equality data: the cell for each clock and the
weak order of its bounded fractional parts. Equality of these data is
reflexive, symmetric and transitive, so it really is an equivalence relation.
Clock \(x\) has \(2M_x+2\) cells. For fixed cells, weak orders on the
nonzero fractional parts can be encoded by a permutation of the clocks
and a cut/no-cut choice between consecutive permutation entries: entries
without a cut belong to one tied block. Include the bounded zero clocks
as a distinguished least block. There are at most \(n!2^n\) encodings;
infeasible/repeated encodings only increase this upper bound. Multiplication
by the choices of cells gives the result. With no clocks there is one
descriptor.

An integer unary threshold at most \(M_x\) distinguishes only which
point/open cell the clock occupies, not its exact fractional value.
Negative thresholds are constant tests on nonnegative values. Thus every
allowed unary atom, and every finite conjunction of them, is constant on
a region. A reset replaces its clock's cell by the zero point, deletes
its old fractional position, and inserts it in the zero block. The order
of all unreset bounded clocks and every over-limit flag stays unchanged.
Hence resets preserve equivalence.

We now prove time matching, including all intermediate times required by
invariants. For each initially bounded clock \(x\), list its nonnegative
integer-hit times
\[
 t_{x,k}=k-v(x),\qquad
 k\in\mathbb Z,\quad v(x)\leq k\leq M_x.
\]
List the corresponding times \(t'_{x,k}=k-w(x)\). The same indices occur
because the floors and zero flags agree. Include time zero as a separate
marked event even when no clock is integral. These two finite lists have
the same weak ordering and ties. To check this, write
\(v(x)=m_x+f_x\), \(v(y)=m_y+f_y\). The difference of event times is
\[
 (k-m_x)-(j-m_y)-(f_x-f_y).
\]
If the integer term is positive it is at least one and the whole expression
is positive because \(-1<f_x-f_y<1\); if negative it is negative.
When it is zero, its sign and vanishing are determined exactly by the
recorded fractional order. The same argument applies to \(w\).
Zero events correspond precisely to integral clocks already at their
listed integer, so their ties with zero agree too.

Map each distinct ordered event time to its counterpart, linearly on
every interval between consecutive events. After the last event use any
positive-slope affine continuation; if zero is the only event, use
\(h(t)=t\). This defines a continuous strictly increasing unbounded
map with \(h(0)=0\), hence a bijection of the nonnegative half-line.
At corresponding events, the same clocks attain the same integer points.
Between events, the same bounded clocks occupy the same integer bands;
those crossing their last point become over-limit in both trajectories.
No bounded clock changes its fractional-order position except when one
of those integer hits occurs. Equivalently, subtracting a common delay
from the difference of two fractional parts leaves their order unchanged
until a wrap at an integer event. Thus every corresponding open interval
has the same fractional weak order. Initially over-limit clocks stay
over-limit forever and contribute no event. This proves
\(v+t\sim_Mw+h(t)\) for every \(t\geq0\), including boundaries.

For a combined edge of delay \(d\), take \(d'=h(d)\). Since \(h\)
maps \([0,d]\) onto \([0,d']\), every intermediate invariant test
transfers by unary stability. The end guard transfers, and reset stability
and destination-invariant stability transfer the final state. Use the same
argument in the reverse direction, or \(h^{-1}\), to obtain bisimulation.
No numerical equality \(d=d'\) was assumed. Strict increase does preserve
whether the delay is zero or positive.
\end{proof}
\end{annexitem}

\begin{annexitem}{annex:p09-timed-region-reachability}{thm:p09-timed-region-reachability}{Projection, path lifting and effective region edges}
\begin{proof}
Every concrete step projects to an edge by the quotient definition, so
every finite run projects to a quotient path. Conversely lift a quotient
path by induction on its length. Start with the actual zero valuation.
Suppose an actual reachable state in the current vertex has been chosen.
The next quotient edge has some representative witnessing its concrete
step. The actual current state is equivalent to that representative;
the bisimulation theorem supplies a step with the same edge/action from
the actual state into the prescribed destination vertex. Continue for the
finitely many edges. The target vertex specifies a location in \(F\),
so the lifted endpoint has the desired property. This induction prevents
the unsupported splicing of incompatible witnesses.

For effectiveness, enumerate all cell assignments and weak fractional
orders and discard inconsistent descriptors. Give each tied nonzero
fractional block a distinct rational value \(i/(b+1)\), in increasing
order for \(b\) blocks. Combine these fractions with the recorded floors,
put point clocks at their integer, and over-limit clocks at \(M_x+1\).
This constructs a rational representative of every feasible descriptor.
For it, enumerate zero, all finitely many integer-hit times, one rational
point in each intervening open interval, and one rational point beyond
the last event. Those samples exhaust all time-successor regions.
Guards are constant on each sampled region, and applying an edge's reset
determines one destination region. Invariant validity along the delay can
either be checked on every preceding event/open cell, or just at its two
endpoints: each conjunct is a unary linear comparison along the delay,
so validity at both endpoints implies validity throughout the segment,
including strict comparisons. Test the destination invariant as well.
The time-matching proof shows that every possible transition from every
representative is represented by this finite enumeration. There are
\(|L|R_M\) or fewer vertices and finitely many edges, so a finite graph
search decides reachability. It makes no claim that quotient cycles meet
infinite-run progress requirements.
\end{proof}
\end{annexitem}

\begin{annexitem}{annex:p09-timed-diagonal-refinement}{thm:p09-timed-diagonal-refinement}{Reset-sensitive diagonal information}
\begin{proof}
Adding \(|\mathcal D|\) truth bits multiplies the descriptor bound by at
most \(2^{|\mathcal D|}\); inconsistent bit combinations are simply empty.
All specified guard/invariant atoms now have the same truth value at
equivalent valuations. Common delay leaves each difference \(x-y\)
unchanged, so the ordinary event-time matching also preserves every added
truth bit for all times. For resets examine a diagonal atom in four cases.
With neither clock reset its truth stays as before. With both reset it is
the constant test \(0\bowtie k\). With just \(y\) reset it becomes the
unary comparison \(x\bowtie k\); with just \(x\) reset it becomes
\(-y\bowtie k\), or the equivalent reversed unary comparison at
\(-k\). The common bound covers \(|k|\), so the ordinary unary region
property settles both last cases. Hence reset stability holds. Transfer
all intermediate invariant tests by the matched continuous map as before;
finite quotient path lifting is unchanged.

A refined descriptor is a finite Boolean combination of rational
difference comparisons: the ordinary cells/fractional order are
difference comparisons after subtracting their fixed integer parts, and
the new bits are atoms or their negations. Convert it into a finite
disjunction of conjunctive difference systems. Negated equalities split
into two strict comparisons; reversed comparisons reverse indices.
The DBM consistency theorem below decides every conjunction and supplies
a rational representative when feasible. This finite splitting is not
an assertion that every refined region is one convex zone. For each
feasible representative, enumerate its integer events as in the ordinary
construction. Diagonal truths remain fixed during delay, so that finite
sample list again suffices for edges. Guard/invariant predicates, resets
and destination descriptors are computable. This supplies the promised
effectiveness. The body example demonstrates why the added bits cannot
be dropped when diagonal guards are present.
\end{proof}
\end{annexitem}

\begin{annexitem}{annex:p09-timed-dbm-closure}{thm:p09-timed-dbm-closure}{Shortest paths with strictness and rational feasibility}
\begin{proof}
Adding comparisons along a path gives its summed comparison; if at least
one premise is strict, its sum is strict. Every closure update therefore
adds an already implied constraint. The new conjunction implies the old
one because bounds are only tightened, giving equality of the solution
sets. The Floyd--Warshall induction says that after processing vertices
\(0,\ldots,k\), each entry is the minimum path bound with intermediate
vertices in that set: such a path either avoids \(k\), or splits there
into its two already considered parts. Initialize diagonals with the
zero-length path. In the absence of a bad cycle a minimum can be taken
on a simple path, since deleting a positive cycle makes its bound tighter
and deleting a zero nonstrict cycle leaves it unchanged. A zero strict
cycle or a negative cycle is bad. Any bad closed walk contains a bad
simple cycle by decomposition into simple cycles. The recurrence therefore
finds a diagonal bound below \((0,\leq)\) whenever such a cycle exists.
Such a bound asserts \(0<0\) or \(0\leq c<0\), impossible for a solution.

It remains to prove the converse, including strict bounds, rather than
assume a shortest-path criterion for reals. Suppose there is no bad cycle.
There are finitely many simple directed cycles. Each zero-total-weight
cycle has no strict edge; every cycle with any strict edge has positive
total constant. Choose a rational \(\varepsilon>0\) so small that
\(N\varepsilon<c\) for every positive simple-cycle constant \(c\),
where \(N=n+1\); choose, for example, half the minimum \(c/N\).
If there are no positive cycles, any positive rational works. Replace
every strict edge constant \(c\) by \(c-\varepsilon\), making the new
system nonstrict. Its simple-cycle constants are nonnegative: each cycle
loses at most \(N\varepsilon\), and a zero cycle loses nothing. Hence
no closed walk is negative.

Add a new super-source with a zero edge to every coordinate and no edges
back into it. Every vertex is reachable and no new cycle is introduced.
Let \(p_i\) be the shortest distance from it to \(i\); the minimum exists
on a simple path and is rational. An edge \(j\to i\) of modified constant
\(c'\) gives \(p_i\leq p_j+c'\). Set \(v_i=p_i-p_0\), so \(v_0=0\).
Every original nonstrict comparison holds, and every original strict one
holds because \(v_i-v_j\leq c-\varepsilon<c\). Included nonnegativity
rows give \(p_0-p_i\leq0\), hence \(v_i\geq0\). This is the required
rational clock valuation. The argument is constructive finite rational
arithmetic. The three-index recurrence uses \((n+1)^3\) updates, proving
the stated bound-operation count without a claim about bit-cost.
\end{proof}
\end{annexitem}

\begin{annexitem}{annex:p09-timed-dbm-inclusion}{thm:p09-timed-dbm-inclusion}{Necessity as well as sufficiency of canonical-bound comparison}
\begin{proof}
If every bound of \(D\) is at least as tight as that of \(E\), every
\(D\)-valuation satisfies \(E\). For necessity suppose
\(D_{ij}\not\preceq_bE_{ij}\). The \(E\) entry cannot be infinity.
Write it as \((c,\triangleleft)\). Add to \(D\) the negation of this
one comparison: the reversed bound \((-c,<)\) when
\(\triangleleft=\leq\), or \((-c,\leq)\) when
\(\triangleleft=<\).

A new bad cycle would have to use this added edge, since \(D\) is
consistent. A cycle using it once consists of that edge and an old path
whose tightest bound is the canonical \(D_{ij}\). If \(D_{ij}\) has
constant larger than \(c\), the cycle constant is positive; if it is
infinity no return path exists. The only remaining possibility is
\(E_{ij}=(c,<)\) and \(D_{ij}=(c,\leq)\). Then the new cycle has
constant zero and no strictness. There cannot be an equally short strict
old path in this last case, since canonical minimum would then have
made \(D_{ij}\) strict. Cycles using the added edge repeatedly split
into such walks and old nonbad cycles, and still are not bad. Therefore
the strengthened system is consistent. The preceding feasibility proof
supplies a valuation satisfying \(D\) and violating the chosen \(E\)
entry, refuting inclusion. Empty-marker cases are immediate from set
inclusion and are handled before this argument.
\end{proof}
\end{annexitem}

\begin{annexitem}{annex:p09-timed-dbm-operations}{thm:p09-timed-dbm-operations}{Projection, simultaneous reset, common delay and intersection}
\begin{proof}
Intersection combines all inequalities from both inputs, which is
entrywise minimum; closure preserves its solution set. The result is
marked empty if the proved cycle test fails.

We prove the projection fact needed for reset. Restrict a canonical
nonempty \(D\) to any set \(K\) of coordinates containing zero. Every
original valuation restricts to a solution of that submatrix. Conversely
take values for \(K\) satisfying it and add a missing coordinate \(r\).
Its upper requirements are \(v_r\mathrel{\triangleleft}v_j+c_{rj}\);
its lower requirements are
\(v_r\mathrel{\triangleright}v_i-c_{ir}\), for \(i,j\in K\), with
strictness reversed in the lower inequality as appropriate. Every lower
comparison \(\triangleright\) here means \(\geq\) for an old nonstrict
bound and \(>\) for an old strict bound; \(\triangleleft\) likewise
means \(\leq\) or \(<\) according to the old bound's flag. Every lower
endpoint is below every upper endpoint by
\[
 D_{ij}\preceq_bD_{ir}\oplus_bD_{rj}
 \quad\text{and}\quad v_i-v_j\text{ satisfies }D_{ij}.
\]
If an endpoint requirement is strict, equality of those endpoints is
excluded by that same strict summed bound. The largest of finitely many
lower endpoints and smallest of finitely many upper endpoints thus give
a nonempty interval, with the correct endpoint flags; infinite missing
bounds cause no obstruction. Choose a point in it. Nonnegativity is
one of its lower requirements. It satisfies the submatrix on
\(K\cup\{r\}\), whose inherited triangle inequalities remain valid.
Repeat for all missing coordinates. This proves exact projection, not
merely inclusion of the projected solutions.

For resetting, project away the old reset coordinates, then insert new
coordinates identically zero. Copying index zero by \(\rho\) does
precisely this: its two bounds against coordinate zero force every new
reset coordinate to be zero, the kept submatrix is unchanged, and every
mixed comparison is its kept-coordinate comparison against zero. The
projection proof supplies compatible old reset values for every new
solution. Hence this matrix describes exactly the reset image. Its
triangle inequalities follow directly from those of \(D\), even when
several indices map to zero. Empty inputs are kept empty separately;
this is why the formula is not applied to an inconsistent raw matrix.

For time-up, every common delay preserves pair differences and can only
weaken the retained lower bounds against zero. Thus the forward image
satisfies the matrix with upper bounds \(D_{i0}\) removed. Conversely
let \(w\) satisfy that relaxed matrix. Seek \(d\geq0\), set
\(v_0=w_0=0\), and put \(v_i=w_i-d\) only for the actual clock
indices \(1\leq i\leq n\). Each old finite upper bound requires
\(d\geq w_i-c_{i0}\), strictly if that upper bound was strict. Each
old lower bound requires \(d\leq w_j+c_{0j}\), strictly if that
lower bound was strict. Canonicality gives
\[
 D_{ij}\preceq_bD_{i0}\oplus_bD_{0j},
\]
so every lower endpoint for \(d\) is below every upper endpoint, with
strict separation whenever either corresponding flag demands it. The
additional lower endpoint zero is compatible with every upper endpoint
because \(w\) already satisfies the retained \(D_{0j}\) bound.
The same finite interval argument supplies \(d\). Included
\(D_{0j}\preceq_b(0,\leq)\) ensures \(v_j\geq0\). For
\(i,j>0\), the differences \(v_i-v_j=w_i-w_j\) still satisfy the
retained clock-to-clock entries. Comparisons with coordinate zero are
instead supplied by the chosen \(d\): the restored upper bounds and
the retained lower bounds were precisely its interval requirements.
Therefore \(w=v+d\) with \(v\in Z(D)\), proving exactness. Strict
endpoints have been respected, not rounded to closed ones.
\end{proof}
\end{annexitem}

\begin{annexitem}{annex:p09-timed-zone-post}{thm:p09-timed-zone-post}{All intermediate invariant tests and symbolic path exactness}
\begin{proof}
A concrete successor lies in the displayed post: its endpoint before
reset is in time-up, satisfies the source invariant and guard, and its
reset satisfies the destination invariant. Conversely any valuation in
the displayed post has witnesses \(v\in Z\) and \(d\geq0\), with
pre-reset endpoint \(w=v+d\) satisfying the source invariant and guard,
and post-reset endpoint in the destination invariant. It remains to
justify the entire delay, not just its end. If \(d>0\),
\(v+t=(1-t/d)v+(t/d)w\) for \(t\in[0,d]\). Each invariant
conjunct is a linear difference comparison. If its two endpoint values
are \(\leq c\), their convex combination is \(\leq c\); if both
are \(<c\), their combination is \(<c\), including endpoints.
Equalities are pairs of weak comparisons. The same reasoning treats
reversed bounds. Since \(Z\) is a subset of the source invariant, both
endpoints satisfy every conjunct. For \(d=0\) the claim is immediate.
Thus this really is an allowed concrete delay/action step.

Start with the singleton zero zone, after checking the initial invariant.
Induct on the chosen finite edge sequence using the exact one-step post
equivalence. Every symbolic endpoint valuation is reached by some
concrete run of that sequence, and every such run has its endpoint in
the symbolic zone. Intersection, time-up and reset closure prove that
the successor remains a zone and supply its exact DBM. No convex-hull
union or finite-state assumption was used. The infinite-zone example
computes \(x=0,y=k\) inductively: the next delay must be exactly one
because its guard is \(x=1\), then only \(x\) is reset. It consequently
does not admit a terminating exhaustive raw-zone search claim.
\end{proof}
\end{annexitem}

\begin{annexitem}{annex:p09-timed-finite-normalization}{thm:p09-timed-finite-normalization}{Saturation commutes with posts at the region level}
\begin{proof}
If \(u\in\alpha(Z)\), some \(v\in Z\) is region equivalent to it.
For any allowed edge step from \(u\), bisimulation supplies the same
edge from \(v\) into an equivalent successor. Conversely, \(Z\subseteq
\alpha(Z)\). Therefore the sets of successor regions of \(Z\) and
\(\alpha(Z)\) are identical, proving the displayed saturation identity.
Guard and invariant constancy also ensures that saturation of a set of
valid states stays within the same location invariant.

The initial zero region is a singleton. Induct on a normalized symbolic
path, using that identity at every edge, to show that its region set
equals saturation of the exact path's endpoint set. Thus any nonempty
target-region set witnesses a concrete finite run, by the finite lifting
proof; conversely every concrete run projects to one. For a single
location there are at most \(2^R\) unions of regions, so the node bound
follows. A more direct algorithm accumulates reachable region entries at
each location and adds the successors of each newly discovered entry.
It stops after at most \(|L|R\) additions. Every addition lifts to a
concrete finite run, and induction on concrete path length proves that
none is omitted.

All required operations are effective. The ordinary cells and fractional
orders are rational difference systems; in the refined case split their
Boolean difference conditions into finitely many conjunctive branches.
The consistency proof supplies rational representatives and decides
intersection with any DBM zone. Event enumeration supplies the finite
region successors. Exact unions are represented by their finite region
IDs, not asserted to have one DBM. Applying the ordinary partition where
diagonal atoms are unstable invalidates the first bisimulation step, as
the displayed diagonal counterexample shows. This proof supplies the
finite normalization used here and no correctness theorem for an
unimplemented practical extrapolation operator.
\end{proof}
\end{annexitem}
